\title{Ten Advances in Mathematics and Theoretical Computer Science}
\begin{document}
\maketitle

\begin{abstract}
We present a collection of results obtained by an internal OpenAI model, spanning mathe- matics and theoretical computer science:
\end{abstract}

\section{Recovered Text}
Ten Advances in Mathematics
and Theoretical Computer Science
OpenAI
Abstract
We present a collection of results obtained by an internal OpenAI model, spanning mathe-
matics and theoretical computer science:
1.
High-dimensional sphere packing. The asymptotic strength of the Cohn–Elkies
linear program is determined exactly. This gives an improved general packing bound
in high dimensions and settles the corresponding Fourier sign-uncertainty problem
asymptotically.
2.
Binary and spherical codes. Classical upper bounds for fixed-distance binary and
spherical codes are improved by exponential factors for all parameters. The spherical
construction also recovers the sphere-packing exponent of Chapter 1.
3.
Non-sofic groups. An explicit non-sofic group is constructed, resolving the question
of whether every countable group admits finite permutation approximations. The
argument uses property-(T) expanders and the binary Leavitt algebra.
4.
Connes’s rigidity conjecture. Infinitely many pairwise nonisomorphic property-
(T) groups are constructed with the same group von Neumann algebra, disproving
Connes’s conjecture and answering a related finite-to-one question of Popa.
5.
Arithmetic circuit complexity. For the permanent, division-free circuits require
Ω(n2 log log n) gates, while formulas require Ω(n4/ log n) leaves.
6.
Quantum parallel repetition. Exponential parallel repetition is proved for every
finite two-player entangled game, extending the classical repetition principle beyond
previously treated special classes of quantum games.
7.
Closest vector problem. A direct reduction from 3SAT gives n1/400-factor hard-
ness for the Euclidean closest vector problem, with related consequences for binary
decoding and other lattice norms.
8.
Ehrhart’s volume conjecture. The sharp bound (n + 1)n/n! is proved in every
dimension for convex bodies whose barycenter is their only interior lattice point.
9.
Multicolor Ramsey numbers. A superexponential lower bound proves Rk(3) =
kΘ(k).
10.
Compactness and degeneracy. Separate bipartite graph constructions disprove
two conjectures in extremal graph theory: the compactness conjecture of Erdős and
Simonovits and a degeneracy conjecture of Erdős.
i
Chapter 1
Exponential Growth Rate of the Cohn–Elkies
Sphere Packing Linear Program
Abstract. We determine the exact exponential decay rate of the Cohn–Elkies
sphere-packing linear program. If ∆d denotes the maximal sphere-packing den-
sity in Rd and LPd denotes the optimal density bound furnished by this program,
then
lim sup
d→∞
∆1/d
d
≤lim
d→∞LP1/d
d
=
r e
2π.
The positive- and negative-eigenvalue Fourier sign-uncertainty radii are both
(1/π + o(1))
√
d.
Contents
1.
Introduction
2.
Fourier-analytic preliminaries
3.
The universal Cohn–Elkies lower bound
4.
The admissible primal upper bound
Appendix A.
Comparison of the sign-uncertainty constants
References
1
1. Introduction
Sphere packing asks how densely congruent balls can fill Euclidean space. For centers sep-
arated by at least 1, let ∆d be the supremum of the upper densities of radius-1/2 balls in Rd.
The Fourier-analytic linear-programming method of Gorbachev and Cohn–Elkies [Gor00, Coh02,
CE03] bounds this density using an auxiliary function and its Fourier transform. Viazovska used
this framework to prove that the E8 lattice gives the optimal packing in dimension eight [Via17];
Cohn, Kumar, Miller, Radchenko, and Viazovska subsequently established the optimality of the
Leech lattice in dimension twenty-four [CKMRV17]. Related Fourier interpolation formulas re-
veal further structure behind these exceptional configurations [RV19, CKMRV22].
By contrast, the behavior of the linear program in high dimensions remained poorly under-
stood. Cohn and Zhao proved that it is always at least as strong as the Kabatianskii–Levenshtein
spherical-code bound [CZ14], but whether it improves the classical high-dimensional sphere-
packing exponent remained open. Motivated by the modular bootstrap [HMR19], Afkhami-
Jeddi, Cohn, Hartman, de Laat, and Tajdini conjectured its high-dimensional rate [AJCHLT20,
§ 3].
We use the Fourier convention and unit-ball volume
bf(ξ) =
Z
Rd f(x)e−2πix·ξ dx,
vd =
πd/2
Γ(d/2 + 1).
(1)
A ball of radius 1/2 has volume vd/2d. For the real Schwartz space S(Rd; R), define
Ad =
n
f ∈S(Rd; R) : bf(0) > 0,
bf ≥0 on Rd,
f ≤0 on \{|x| ≥1\}
o
,
(2)
and the associated linear-programming bound by
LPd = vd
2d inf
f∈Ad
f(0)
bf(0)
.
(3)
Fourier inversion gives f(0) > 0. The Gorbachev–Cohn–Elkies bound [Gor00, CE03] assigns
every f ∈Ad the density upper bound vdf(0)/(2d bf(0)), so (3) yields
∆d ≤LPd.
(4)
Our main theorem proves the exponential rate conjectured in [AJCHLT20, Conjs. 3.1–3.2].
Theorem 1.1. As d →∞, LP1/d
d
−→
p
e/(2π).
By (4), the theorem gives ∆d ≤LPd = 2−(α∗+o(1))d, where α∗= 1
2 log2(2π/e) = 0.6044 . . ..
This is the first improvement since 1978 to the general sphere-packing exponent. The classi-
cal Kabatianskii–Levenshtein exponent was 0.59905576 . . . [KL78, Lev79]; subsequent spherical-
code refinements had improved only lower-order factors [CZ14, SZ24, Z24]. The matching lower
bound shows that no Cohn–Elkies auxiliary function can improve this exponent.
Our companion paper in Chapter 2 recovers the upper-bound direction of Theorem 1.1 as a
small-distance limit of spherical-code bounds.
The packing sign conditions are related to the Bourgain–Clozel–Kahane uncertainty princi-
ple for eventually nonnegative Fourier eigenfunctions [BCK10, GOSS17, GOSR21]. Cohn and
Gonçalves [CG19] introduced its anti-self-Fourier counterpart and connected it with sphere
packing. If g ∈L1(Rd; R) satisfies bg = ςg for ς ∈\{−1, +1\}, then bg ∈L1 and Fourier inversion
supplies a continuous representative of g. Using this representative for all pointwise values and
sign conditions, define
r(g) = inf\{R ≥0 : g(x) ≥0 for |x| ≥R\},
(5)
Aς(d) = inf\{r(g) : 0 ̸= g ∈L1(Rd; R), bg = ςg, g(0) = 0\}.
(6)
Here r(g) = ∞when no such radius exists. The signs +1 and −1 give the original and comple-
mentary uncertainty problems, respectively.
The negative-eigenvalue problem also arises in the spinless modular bootstrap, where mod-
ular S-antisymmetry produces anti-self-Fourier test functions whose eventual sign controls the
2
spectral gap [HMR19, AJCHLT20].
General lower bounds for sign uncertainty and limita-
tions of Gaussian–polynomial test functions of sublinear degree appear in [Edw25] and [CDG24,
Thm. 1.2], respectively. The next theorem proves the asymptotic equality conjectured in [CG19,
Conj. 1.5] and [AJCHLT20, (3.5)], with the common value predicted in [AJCHLT20, (3.3)]. Al-
though these asymptotics coincide, Appendix A observes that A+(d) < A−(d) for all d.
Theorem 1.2. The Fourier-eigenfunction sign-uncertainty constants satisfy
lim
d→∞
A+(d)
√
d
= lim
d→∞
A−(d)
√
d
= 1
π.
The proof reduces both problems to the last sign changes of Fourier eigenfunctions. In § 3,
Proposition 3.1 shows that a radial Schwartz function g satisfying bg = ±g and g(0) = 0 has
exponentially little L1 mass inside B(0, c
√
d) whenever c < 1/π. For an admissible packing
function F, let a = ( bF(0)/F(0))1/d and h(x) = F(ax). Then h(0) = bh(0), while the nonzero
function g = bh−h is anti-self-Fourier, vanishes at the origin, and is nonnegative outside B(0, 1/a).
Since
R g = 0, its negative part has mass ∥g∥1/2, all of which lies in this ball. The mass estimate
therefore forces 1/a ≥(1/π −o(1))
√
d. In § 4, Theorem 4.1 modifies the Mellin transform of
a Gaussian to construct a Fourier pair and a self-Fourier function whose required exterior sign
conditions begin at (1/π + o(1))
√
d. Stirling’s formula converts the matching radius bounds
into the packing exponent
p
e/(2π).
2. Fourier-analytic preliminaries
We reduce the packing and sign-uncertainty problems to radial functions and establish the
Mellin–Fourier identities used in both bounds. We repeatedly use the gamma recurrence and
reflection formulas [DLMF, § 5.5], together with their consequences for real b ̸= 0:
Γ(z + 1) = zΓ(z),
Γ(z)Γ(1 −z) =
π
sin(πz),
|Γ(ib)|2 =
π
b sinh(πb),
|Γ(1/2 + ib)|2 =
π
cosh(πb).
(7)
2.1. Radial reduction. Throughout, Fourier transforms use the convention (1), and we denote
a radial function and its one-variable profile by the same symbol. With normalized Haar measure
on O(d), define the rotational average
Rf(x) =
Z
O(d)
f(Ux) dU.
Write Srad(Rd; R) for the real Schwartz functions depending only on |x|, and L1
rad(Rd; R) for the
radial real integrable functions. Set Arad
d
= Ad ∩Srad(Rd; R). For every continuous integrable
f, rotational invariance gives
d
Rf = R bf,
(Rf)(0) = f(0),
d
Rf(0) = bf(0).
Consequently, rotational averaging preserves every sign condition in (2), so the infimum in (3)
is unchanged when Ad is replaced by Arad
d
[CE03, CG19]. If bg = ςg, then
d
Rg = ςRg,
r(Rg) ≤r(g).
Moreover, Rg ̸= 0 whenever g ̸= 0 is eventually nonnegative. Indeed, if Rg = 0, the nonnegative
values of g outside a ball have zero rotational average, so g vanishes outside that ball. Its Fourier
transform is then entire, and bg = ςg makes this entire function vanish on the same exterior region,
forcing g = 0.
We also need approximation by Schwartz functions preserving the Fourier eigenvalue and
vanishing at the origin. Suppose that g ∈L1
rad(Rd; R) satisfies bg = ςg and g(0) = 0, where
ς ∈\{−1, +1\}. For n ≥1, set
κn(x) = nde−πn2|x|2,
ηn(x) = e−π|x|2/n2,
qn = ηn(g ∗κn).
3
Gaussian mollification gives qn ∈Srad(Rd; R), qn →g, and bqn = κn ∗(ηnbg) →bg in L1. Moreover,
qn(0) →g(0) = 0 and bqn(0) =
R qn →bg(0) = 0. Define
pn = qn + ς bqn
2
,
ψ+(x) = e−π|x|2,
ψ−(x) =

|x|2 −d
4π

e−π|x|2.
Since bpn = ςpn, bψς = ςψς, and ψς(0) ̸= 0, the correction
gn = pn −pn(0)
ψς(0)ψς
satisfies gn ∈Srad(Rd; R), bgn = ςgn, gn(0) = 0, and gn →g in L1.
2.2. The radial Mellin transform. For a radial Schwartz function g, write r = |x| and ρ = |ξ|
for the spatial and frequency radii, respectively, and set
λ = d
2,
Sd = 2πd/2
Γ(d/2).
Here Sd is the area of the unit sphere, and polar integration gives
Z
Rd g(x) dx = Sd
Z ∞
0
g(r)rd−1 dr.
For ρ > 0, the Fourier transform has the Hankel representation
bg(ρ) = 2πρ1−d/2
Z ∞
0
g(r)Jd/2−1(2πrρ)rd/2 dr.
Because its Bessel kernel depends on rρ, the radial Fourier transform becomes particularly
simple after a Mellin transform: it reflects the Mellin variable and multiplies by an explicit
gamma factor. For Re z > 0, t ∈R, and r > 0, the Mellin transform, its restriction to the
critical line, and Mellin inversion [PK01] are
Mg(z) =
Z ∞
0
g(r)rz−1 dr,
Xg(t) = Mg(λ −it),
g(r) = r−λ
2π
Z
R
Xg(t)rit dt.
(8)
In particular, the radial integration normalization is
bg(0) =
Z
Rd g(x) dx = SdMg(d).
In logarithmic radius v = log r, set Φg(v) = eλvg(ev). Smoothness at r = 0 gives exponential
decay of Φg and all its derivatives as v →−∞, while the Schwartz decay of g gives rapid decay
as v →+∞. Thus Φg ∈S(R; R), and ordinary one-dimensional Fourier inversion gives
Xg(t) =
Z
R
Φg(v)e−itv dv,
Φg(v) = 1
2π
Z
R
Xg(t)eitv dt.
The standard Mellin transform of the Hankel kernel gives, for 0 < Re z < d, the Mellin–Hankel
functional equation [SW71, Ch. IV], [DLMF, (10.22.43)]:
Mbg(z) = πλ−z
Γ(z/2)
Γ((d −z)/2)Mg(d −z).
(9)
If g(0) = bg(0) = 0, smooth radiality gives g(r), bg(r) = O(r2) at the origin.
Both Mellin
transforms therefore extend holomorphically to Re z > −2, and (9) continues to Re z = 0:
Mg(d) = S−1
d bg(0) = 0,
Mg(d −z) = −zM ′
g(d) + O(z2),
Γ(z/2) = 2
z + O(1),
so the apparent gamma pole at z = 0 cancels. The line Re z = d/2 is fixed by the reflection
z 7→d −z. On this line the Fourier transform acts, for t ∈R, by
Xbg(t) = mλ(t)Xg(−t),
mλ(t) = πit Γ((λ −it)/2)
Γ((λ + it)/2).
(10)
4
For t ∈R, its multiplier satisfies
mλ(−t) = mλ(t) = mλ(t)−1,
|mλ(t)| = 1.
3. The universal Cohn–Elkies lower bound
By the radial reduction in § 2.1, assume that the admissible packing function F is radial. Let
a = ( bF(0)/F(0))1/d and define h(x) = F(ax). Then h(0) = bh(0), and g = bh −h is anti-self-
Fourier, vanishes at the origin, and is nonnegative for |x| ≥1/a; see (29)–(30). The proof of
Theorem 3.8 verifies that g ̸= 0. Since
R g = 0, its negative part has mass ∥g∥1/2, all of which
lies in B(0, 1/a).
Proposition 3.1 shows that every radial Schwartz Fourier eigenfunction vanishing at the origin,
of either eigenvalue, has exponentially little L1 mass in B(0, c
√
d) when c < 1/π. Thus the
negative half-mass of g cannot fit inside this ball, forcing 1/a ≥(1/π −o(1))
√
d.
The obstruction comes from the Mellin–Fourier identity (9). Lemma 3.2 bounds a normalized
Mellin transform Z on the strip | Im t| < d/2: total L1 mass controls the upper boundary,
while the Fourier functional equation controls the lower boundary. Writing λ = d/2, Poisson
interpolation gives
log |Z(s + iσλ)| ≤Hσ(s) ≤λMσ
log(2πc2) + Jσ
 + Oσ(log λ),
Mσ = 1 −σ
2
,
by Lemma 3.3. The sharp constant enters through Lemma 3.4: Jσ →log(π/2) as σ ↑1, so the
parenthesized rate tends to log(π2c2). It is negative exactly when c < 1/π. An all-frequency
bound and shifted Mellin inversion, in Lemmas 3.5 and 3.6, turn this negativity into the interior-
mass estimate; Stirling’s formula then yields the packing lower bound.
3.1. The Mellin-strip obstruction.
Proposition 3.1. For every 0 < c < 1/π, there exist Cc, γc > 0 and d0(c) ∈N such that, for
every d ≥d0(c), every ς ∈\{−1, +1\}, and every nonzero g ∈Srad(Rd; R) satisfying bg = ςg and
g(0) = 0, one has
Z
|x|<c
√
d
|g(x)| dx ≤Cce−γcd∥g∥1.
(11)
Fix 0 < c < 1/π, d ∈N, λ = d/2, R = c
√
d, and a nonzero g ∈Srad(Rd; R) with bg = ςg,
ς ∈\{−1, +1\}, and g(0) = 0.
Let Sd = 2πλ/Γ(λ) be the area of the unit sphere.
In the
logarithmic coordinate r = Rev, normalize the Mellin inversion formula (8) by setting
φ(v) =
Sd
∥g∥1
(Rev)dg(Rev),
Z(t) =
Sd
∥g∥1
Rλ+itXg(t),
(12)
∥φ∥1 = 1,
Z
R
φ = 0,
Z 0
−∞
|φ(v)| dv =
1
∥g∥1
Z
|x|<R
|g(x)| dx.
(13)
The second line follows from polar integration and bg(0) = ςg(0) = 0.
Lemma 3.2. For every −1 < σ < 1, the function Z is bounded and holomorphic on a neighbor-
hood of the strip |Im t| ≤λ. Its boundary values satisfy |Z(y+iλ)| ≤1 and log|Z(y−iλ)| ≤hλ(y)
for y ̸= 0, where the lower-boundary majorant is
hλ(y) = λ log(πR2) + log |Γ(−iy/2)| −log |Γ(λ + iy/2)|.
(14)
Writing θ = π(1 + σ)/2, define
Pσ(T) =
sin θ
4(cosh(πT/2) −cos θ),
Mσ =
Z
R
Pσ = 1 −σ
2
.
(15)
Then
|Z(s + iσλ)| ≤exp
Z
R
Pσ(T)hλ(s −λT) dT

(s ∈R).
5
Proof. Since g(0) = bg(0) = 0, the holomorphic continuation and pole cancellation following
(9) show that Z is holomorphic whenever Im t > −λ −2. In particular, it is holomorphic on
a neighborhood of the closed strip. Its upper-boundary values follow from (12), whereas the
Mellin–Fourier functional equation (9) gives the lower-boundary values:
Z(y + iλ) =
Z
R
φ(v)e−iyv dv,
|Z(y + iλ)| ≤1,
(16)
Z(y −iλ) = ς(πR2)λ+iy Γ(−iy/2)
Γ(λ + iy/2)Z(−y + iλ).
(17)
Taking absolute values and using (16) gives log |Z(y −iλ)| ≤hλ(y), independently of ς. At
y = 0, the zero Z(iλ) = 0 from (13) cancels the gamma pole in (17). Hence Z(y −iλ) remains
bounded there, although
hλ(y) = −log |y| + Oλ(1)
(y →0).
For t0 = s+iσλ, the conformal map t 7→exp(π(t+iλ)/(2λ)) and the upper-half-plane Poisson
formula [Ahl79] give the lower-edge harmonic measure λ−1Pσ((s −y)/λ) dy, with Pσ and Mσ
as in (15). The complementary upper-edge measure is (1 + σ)/2, and its boundary bound (16)
contributes nothing.
The logarithmic singularity of hλ requires a bounded truncation before applying the Poisson
principle. The actual lower-boundary values satisfy
sup
y∈R
|Z(y −iλ)| ≤SdRd
∥g∥1
Z ∞
0
|g(r)|
r
dr < ∞.
Indeed, g(r) = O(r2) at zero and is Schwartz at infinity. More generally, for −λ ≤η ≤λ, (12)
gives
|Z(s + iη)| ≤SdRλ−η
∥g∥1
Z ∞
0
|g(r)|rλ+η−1 dr.
Splitting at r = 1 bounds this expression uniformly in s and η, so Z is bounded on the closed
strip. Choose D > max\{0, supy log |Z(y −iλ)|\}. The strip Poisson principle [Ahl79], applied to
log |Z| with lower-boundary majorant min\{hλ, D\} and upper-boundary majorant 0, gives
log |Z(t0)| ≤
Z
R
1
λPσ
s −y
λ

min\{hλ(y), D\} dy.
The kernel in (15) decays exponentially, the singularity of hλ is locally integrable, and Stirling’s
formula gives hλ(y) = −λ log |y| + Oλ(1) at infinity [DLMF, § 5.11]. Dominated convergence
therefore permits D →∞, and the substitution T = (s −y)/λ yields
|Z(s + iσλ)| ≤eHσ(s),
Hσ(s) =
Z
R
Pσ(T)hλ(s −λT) dT.
□
(18)
Lemma 3.3. For every −1 < σ < 1, there exists Cσ > 0, independent of d, c, and g, such that
Z
R
Pσ(T)
hλ(λT) −λ

log(2πc2) −
Z 1
0
log
q
x2 + T 2/4 dx
 dT ≤Cσ log(2 + λ).
(19)
Define
Jσ = −1
Mσ
Z
R
Pσ(T)
Z 1
0
log
q
x2 + T 2/4 dx dT.
Then, for every s ∈R,
Hσ(s) ≤Hσ(0) = λMσ
log(2πc2) + Jσ
 + Oσ(log(2 + λ)).
(20)
Proof. To estimate the lower-boundary majorant (14) with R = c
√
2λ, put
fT (x) = log
q
x2 + T 2/4,
b = λT
2 .
6
If d = 2n, so that λ = n, the gamma recurrence in (7) and |Γ(−ib)| = |Γ(ib)| give, for T ̸= 0,
hn(nT) = n log(2πc2) −
n−1
X
k=0
fT (k/n).
Monotonicity of fT bounds the Riemann-sum error by fT (1) −fT (0); thus
0 ≤hn(nT) −n

log(2πc2) −
Z 1
0
fT (x) dx

≤1
2 log

1 + 4
T 2

.
If d = 2n + 1, so that λ = n + 1
2, the gamma identities (7) instead give
hλ(λT) = λ log(2πc2) −
n−1
X
k=0
fT
k + 1/2
λ

+ 1
2 log λ + 1
2 log coth(π|b|)
|b|
.
The midpoint Riemann-sum error on [0, n/λ] is at most fT (n/λ)−fT (0); the remaining interval
[n/λ, 1] has length 1/(2λ). Together, these contributions are bounded by
C

1 + log(2 + |T|) + log(2 + |T|−1)

.
Adding C log(2 + λ) also bounds the gamma endpoint correction 1
2 log λ + 1
2 log(coth(π|b|)/|b|).
Since (15) gives Pσ(T) ≪σ e−π|T|/2, and log(2+|T|−1) is locally integrable, integrating the even-
and odd-dimensional bounds proves (19).
It remains to identify the maximum of the Poisson extension Hσ. Both hλ and Pσ are even,
and the digamma series [DLMF, § 5.7] gives, for y > 0,
h′
λ(y) = 1
2 (Im ψ(λ + iy/2) −Im ψ(iy/2)) < 0,
Im ψ(a + ib) =
∞
X
k=0
b
(k + a)2 + b2 .
Here ψ = Γ′/Γ. The explicit kernel (15) is likewise decreasing on (0, ∞). For N > 0, the function
qN = (hλ + N)+ is nonnegative, integrable, even, and decreasing there. The convolution of two
such functions is largest at zero. Subtracting NMσ and letting N →∞, using the exponential
decay of Pσ, therefore gives Hσ(s) ≤Hσ(0). Evaluating at zero with (19) proves (20).
□
Lemma 3.4. For Jσ defined in Lemma 3.3,
lim
σ↑1 Jσ = log π
2 .
Consequently, for every 0 < c < 1/π, there exists σ = σ(c) ∈(−1, 1) such that log(2πc2) + Jσ <
0.
Proof. With T = 2u, divide the lower-edge harmonic measure Pσ(T) dT from (15) by its total
mass Mσ. The corresponding probability density in u is
pσ(u) = 2Pσ(2u)
Mσ
=
sin θ
(1 −σ)(cosh(πu) −cos θ).
As σ ↑1, sin θ/(1−σ) →π/2. The densities pσ are uniformly bounded by Ce−π|u| and therefore
converge in L1(R) to
p(u) = π
4 sech2
πu
2

.
The characteristic function of this density is
Z
R
p(u)eitu du =
t
sinh t,
t ∈R,
(21)
interpreted as 1 at t = 0. For t > 0, a contour shift by 2i across the double pole at i gives
(1 −e−2t)
R p(u)eitu du = 2te−t; evenness handles t < 0. Define I(x) =
R
R p(u) log
√
x2 + u2 du.
For x > 0, the Laplace representation of x/(x2 + u2), (21), and the trigamma integral [DLMF,
§ 5.9] yield
I′(x) =
Z ∞
0
e−xt
t
sinh t dt = 1
2ψ′
x + 1
2

.
7
Since both I(x) and ψ((x + 1)/2) + log 2 equal log x + o(1) at infinity [DLMF, § 5.11], their
integration constants agree. Local integrability at u = 0 extends the identity to x = 0:
Z
R
p(u) log
p
x2 + u2 du = ψ
x + 1
2

+ log 2
(x ≥0).
(22)
The uniform exponential bound on pσ and local integrability of log |u| justify dominated con-
vergence in Jσ. Since
Z 1
0
ψ
x + 1
2

dx = 2 log Γ(1)
Γ(1/2) = −log π,
(22) gives the exact threshold
lim
σ↑1 Jσ = log π
2 ,
log(2πc2) + log π
2 = log(π2c2).
If c < 1/π, choose σ = σ(c) < 1 sufficiently close to 1 that
δc = −
log(2πc2) + Jσ(c)
 > 0.
□
(23)
Fix σ = σ(c) and δc > 0 as in (23). We first bound the Mellin transform Z on the horizontal
line Im t = σλ, and then use that bound to control the mass of g inside B(0, c
√
d).
Lemma 3.5. There exist γc, Cc, Bc > 0, depending only on c, such that, for every sufficiently
large d,
Hσ(s) ≤−γcλ
(s ∈R),
Z
R
|Z(s + iσλ)| ds ≤Ccλe−γcλ.
Moreover, after increasing Bc if necessary,
Hσ(λS) ≤−Mσλ
2
log |S|
Cc
(|S| ≥Bc).
Proof. The maximum estimate for Hσ in (20), together with log(2πc2) + Jσ = −δc, gives
Hσ(s) ≤−λMσδc + Oσ(log λ) ≤−γcλ
(s ∈R)
for all sufficiently large d, with γc > 0 independent of s, g, and the Fourier eigenvalue ς.
To control all Mellin frequencies, apply the gamma identities (7) to (14). In both parities,
hλ(λU) ≤λ log 4πc2
|U| + Eλ(U)
(U ̸= 0),
where
Eλ(U) =



0,
λ ∈N,
1
2 log coth
 πλ|U|
2

,
λ ∈N + 1
2.
Each gamma-recurrence factor has modulus at least |λU|/2; in odd dimension, the remaining
ratio at b = λU/2 contributes Eλ(U). Expanding log coth x in its convergent odd-exponential
series gives, in odd dimension,
Z
R
Eλ(U) dU = 2
πλ
Z ∞
0
log coth x dx = π
4λ.
Thus in both parities the Pσ-convolution of Eλ is at most π∥Pσ∥∞/(4λ). Consequently, (18)
gives
Hσ(λS) ≤λMσ log(4πc2) −λ
Z
R
Pσ(T) log |S −T| dT + Oσ(λ−1).
Split the logarithmic integral at |T| = |S|/2. The exponential decay of (15) gives mass Mσ +
Oσ(e−π|S|/4) on |T| ≤|S|/2; the possible negative contribution from |S −T| < 1 is Oσ(e−π|S|/2),
by local integrability of log |S −T|. Hence, for some Bc, C′
c > 0,
Z
R
Pσ(T) log |S −T| dT ≥Mσ
2 log |S| −C′
c
(|S| ≥Bc).
8
After increasing C′
c, the estimates for Hσ on the whole line and at large frequencies become
Hσ(s) ≤−γcλ
(s ∈R),
(24)
Hσ(λS) ≤−Mσλ
2
log |S|
C′c
(|S| ≥Bc).
(25)
Choose B > max\{Bc, C′
c\} and q = Mσλ/2 > 1. The two bounds in (24)–(25) yield
Z
|s|≤Bλ
eHσ(s) ds ≤2Bλe−γcλ,
Z
|s|>Bλ
eHσ(s) ds ≤2λC′
c
q −1
 B
C′c
1−q
.
The second bound decays at rate (Mσ/2) log(B/C′
c) > 0.
Decreasing γc if necessary, and
applying (18), we obtain
Z
R
|Z(s + iσλ)| ds ≤Ccλe−γcλ.
□
(26)
The global L1 estimate (26) for Z on Im t = σλ now controls φ(v) for v < 0, corresponding
by (13) to |x| < c
√
d.
Lemma 3.6. For every 0 < c < 1/π, there exist Cc, γc > 0 and d0(c) ∈N, independent of g
and ς, such that
Z 0
−∞
|φ(v)| dv ≤Cce−γcd
(d ≥d0(c)).
Proof. Let G(v) = e(σ−1)λvφ(v). The normalized profile φ from (12) satisfies
Z
R
|G(v)| dv = SdR(1−σ)λ
∥g∥1
Z ∞
0
|g(r)|r(1+σ)λ−1 dr < ∞.
Thus G ∈L1(R), and (12) identifies its angular-frequency Fourier transform
R
R G(v)e−isv dv
with Z(s+iσλ). The integrability of this transform follows from (26), so Fourier inversion gives
Z(s + iσλ) =
Z
R
e(σ−1)λvφ(v)e−isv dv,
φ(v) = e(1−σ)λv
2π
Z
R
Z(s + iσλ)eisv ds.
Taking absolute values and integrating over v < 0 contributes
R 0
−∞e(1−σ)λv dv = ((1 −σ)λ)−1.
Combining this with (26) and λ = d/2, and decreasing γc if necessary, gives
Z 0
−∞
|φ(v)| dv ≤
1
2π(1 −σ)λ
Z
R
|Z(s + iσλ)| ds ≤Cce−γcd.
□
(27)
Proof of Proposition 3.1. By (13), the left side of (27) is precisely ∥g∥−1
1
R
|x|<c
√
d |g(x)| dx. Thus
(27) proves (11), uniformly in g and its Fourier eigenvalue ς.
□
3.2. The packing lower bound.
Proposition 3.7. For every 0 < c < 1/π, there exists d0(c) ∈N such that, for every d ≥d0(c)
and ς ∈\{−1, +1\}, no nonzero g ∈L1(Rd; R) satisfies bg = ςg, g(0) = 0, and g(x) ≥0 for
|x| ≥c
√
d. Here g denotes its continuous Fourier-inversion representative.
Proof. First, suppose that g is a radial Schwartz eigenfunction. Since
R g = bg(0) = ςg(0) = 0,
its negative part g−= max\{−g, 0\} has integral ∥g∥1/2. The assumption g(x) ≥0 for |x| ≥c
√
d
forces g−to vanish outside B(0, c
√
d), contradicting (11) once Cce−γcd < 1/2.
For general g ∈L1, the radial reduction in § 2.1 shows that h = Rg ̸= 0 and has the
same Fourier eigenvalue, origin value, and exterior sign. The approximation constructed there
gives radial Schwartz eigenfunctions hn →h in L1 with bhn = ςhn and hn(0) = 0. Writing
(hn)−= max\{−hn, 0\} and R = c
√
d, the exterior nonnegativity of h gives
Z
Rd(hn)−≤
Z
|x|<R
|hn(x)| dx +
Z
|x|≥R
|hn(x) −h(x)| dx.
9
Applying (11) to hn therefore implies
1
2∥hn∥1 =
Z
Rd(hn)−≤Cce−γcd∥hn∥1 + ∥hn −h∥1.
Letting n →∞gives ∥h∥1/2 ≤Cce−γcd∥h∥1, which contradicts h ̸= 0 for all sufficiently large
d.
□
Theorem 3.8. There is a sequence ϵd →0 such that, for every sufficiently large d and every
F ∈Ad,
F(0)
bF(0)
≥2d
vd
r e
2π −ϵd
d
.
(28)
Proof. By radial reduction, replace F with its rotational average, which preserves F(0), bF(0),
and admissibility. Define a = ( bF(0)/F(0))1/d > 0, h(x) = F(ax), and g = bh −h. The Fourier
convention (1) and the admissibility conditions (2) give
bh(ξ) = a−d bF(ξ/a),
h(0) = bh(0) = F(0),
(29)
bg = −g,
g(0) = 0,
g(x) = a−d bF(x/a) −F(ax) ≥0
(|x| ≥1/a).(30)
Moreover, g ̸= 0: otherwise h = bh ≥0, while h(x) = F(ax) ≤0 for |x| ≥1/a. Thus h would be
a nonzero compactly supported self-Fourier function, contradicting Fourier analyticity. Hence
Proposition 3.7 gives 1/a > c
√
d, uniformly in F, for each fixed c < 1/π and all sufficiently
large d. Consequently,
lim inf
d→∞
1
√
d
inf
F∈Ad
F(0)
bF(0)
!1/d
≥c
(0 < c < 1/π).
Taking the supremum over c < 1/π yields
inf
F∈Ad
F(0)
bF(0)
!1/d
≥
 1
π −o(1)
 √
d.
(31)
Combining (31) with Stirling’s formula [DLMF, § 5.11], v1/d
d
= (1+o(1))
p
2πe/d gives (28), for
a sequence ϵd →0 independent of F.
□
4. The admissible primal upper bound
Section 3 established
min


inf
F∈Ad
F(0)
bF(0)
!1/d
, A−(d), A+(d)


≥
 1
π −o(1)
 √
d.
We explain how the upper bounds in the introduction reduce to constructing functions whose
sign changes occur at the same radius. Suppose that real radial Schwartz functions f−, f+ satisfy
bf−= f+ > 0,
f−(0) = f+(0),
f−(x) < 0
(|x| ≥R).
For F(x) = f−(Rx), Fourier scaling gives
bF(ξ) = R−df+(ξ/R) > 0,
F(x) < 0
(|x| ≥1),
F(0)
bF(0)
= Rd.
Consequently F ∈Ad and LPd ≤vd(R/2)d. The difference g−= f+ −f−is anti-self-Fourier,
vanishes at the origin, and is positive for |x| ≥R, so it also proves A−(d) ≤R.
A self-
Fourier function f0 satisfying f0(0) = 0 and f0(x) > 0 for |x| ≥R similarly proves A+(d) ≤R.
Thus all the upper bounds reduce to producing one Fourier pair, one self-Fourier function, and
R = (1/π + o(1))
√
d.
10
Theorem 4.1. There is ϵ0 > 0 such that, for every fixed 0 < ϵ < ϵ0 and every sufficiently large
dimension d, there exist real radial Schwartz functions f−, f+, f0 and a radius Rϵ,d > 0 such
that f+ > 0 everywhere, f−< 0 < f0 whenever |x| ≥Rϵ,d, and
bf−= f+,
bf0 = f0,
f−(0) = f+(0) > 0,
f0(0) = 0.
Moreover,
lim
ϵ↓0 lim
d→∞
Rϵ,d
√
d
= 1
π.
4.1. Outline of the construction. We need a Fourier pair f−, f+ = bf−and a self-Fourier
function f0, with f+ > 0 everywhere and f−< 0 < f0 outside a common radius. In Mellin
coordinates, (10) reduces these Fourier symmetries to reflection of the frequency. The Gaussian
gG(r) = 2πλ/2e−πr2,
EG
λ (t) = πit/2Γ
λ −it
2

,
λ = d
2,
already obeys mλ(t)EG
λ (−t) = EG
λ (t).
Multiplication by an even factor therefore preserves
Fourier symmetry. We choose Eλ(t) = EG
λ (t)eλhϵ(t/λ), where the even perturbation hϵ in (36) is
determined by a signed density w.
The variable a in that density parametrizes radial dilations: the Mellin multiplier cos(at/λ)−1
corresponds to
g(r) 7−→eag(rea/λ) + e−ag(re−a/λ)
2
−g(r).
A shell is an interval of these dilation parameters supporting one component of w. Our negative
component ws moves the sign radius inward, and a positive component wB, supported on much
larger dilation parameters, restores decay at every nonzero Mellin frequency.
We multiply the common envelope by polynomials P−, P+, P0. Reflection exchanges P−and
P+ and fixes P0, giving the desired Fourier symmetries. The first gamma pole, at normalized
frequency ζ = −i, controls the value at the origin: the polynomial P0 cancels this pole, while
P−and P+ retain equal positive residues. On the imaginary axis, the polynomials are chosen
so that P+(iu) > 0 for u > −1, whereas P−(iu) < 0 and P0(iu) > 0 just above u = 1; see (39).
Thus the essential contour is t = λ(T + iu) with u ≈1.
On this contour, the radius r = ev(u) for which the Mellin integrand is stationary at T = 0
satisfies (44). For fixed u > −1, the digamma asymptotic gives
ev(u)
√
d
−→
r1 + u
4π
exp
Z ∞
0
w(a)a sinh(ua) da

.
Consequently, negative w decreases the logarithm of the stationary radius when u > 0. However,
it also reduces the decay of the Mellin integrand away from T = 0. At u = 1, after dividing
the damping by λ, the limiting gamma contribution has density e−2a/(2a2), while the pertur-
bation contributes w(a) cosh a. A sufficient pointwise condition for nonnegative total damping
is therefore
|w(a)| cosh a ≤e−2a
2a2 .
Within this pointwise constraint, the greatest inward displacement is achieved formally by
w∗(a) = −
e−2a
2a2 cosh a,
Z ∞
0
w∗(a)a sinh a da = −1
2 log π
2 ,
(32)
where the integral evaluation follows from the gamma duplication formula [DLMF, § 5.5]. The
Gaussian stationary radius at u = 1 is (2π)−1/2√
d, so the displacement in (32) gives
1
√
2π exp

−1
2 log π
2

= 1
π.
(33)
Although its Mellin perturbation is well defined, the formal density w∗has infinite mass at
zero and saturates the gamma damping. Its Mellin factor therefore lacks the strict decay needed
for a Schwartz function. Instead, we truncate it to a bounded interval and taper it slightly,
obtaining a negative shell ws that preserves the displacement up to O(ϵ) while leaving a definite
11
damping margin. This margin suffices near u = 1, but not on contours with arbitrarily large u.
A second shell wB, supported on [B, B + 1], is chosen to have negligible effect on the stationary
radius at u = u0 while dominating the negative shell at every frequency when u ≥U > u0.
Using an interval rather than a single dilation avoids frequencies at which its damping would
vanish.
The remaining proof has two analytic parts. Uniform saddle estimates show that the Mellin
integral at r = ev(u) has the sign of Pj(iu), which gives the exterior signs and positivity of f+
for r ≥r∗. A downward contour shift then expresses f+ on [0, r∗] as a perturbed exponential
series, proving positivity at the remaining radii. The negative-shell displacement in (32) finally
yields the sharp radius (33).
4.2. The Mellin ansatz. We turn the ideal negative density w∗in (32) into two compactly
supported shells. The negative shell ws retains its logarithmic-radius displacement up to O(ϵ),
while the positive shell wB is negligible at u = u0 and dominates the damping for u ≥U.
Put λ = d/2, and throughout the construction fix ϵ before taking d →∞. Constants in Oϵ(·)
and ≪ϵ may depend on ϵ, but are independent of d and of the saddle parameter whenever a
bound is stated uniformly in that parameter; unadorned implicit constants are absolute.
Introduce cutoffs 0 < a0 < A < B, a positive-shell amplitude Q > 0, and saddle parameters
u0 = 1 + ϵ
4,
U = 1 + ϵ
2,
C0 = A + a−1
0 .
As ϵ ↓0, the cutoffs should satisfy
a0 = o(ϵ),
e−2A
A
= o(ϵ),
ϵA = o(1),
A = o(B).
The first two conditions make the omitted portions of the ideal saddle displacement negligible,
and the third permits an O(ϵ(1 + a)) taper on the negative shell. For the positive shell, of
amplitude Q, the required separation is
BQe(u0−1)B = o(1),
C0Q−1e−(U−1)(B−A) = o(1).
The bound on BQe(u0−1)B keeps the positive shell from moving the target saddle. The bound
on C0Q−1e−(U−1)(B−A) makes that shell dominate the negative shell once u ≥U.
Because
u0 −1 < U −1, both conditions are compatible: writing Q = e−qϵB, choose u0 −1 < qϵ < U −1
with enough separation that ϵB dominates log B + log C0.
One convenient realization of all these requirements is
a0 = ϵ2,
A = log(1/ϵ),
B = ϵ−3,
qϵ = (u0 −1) + (U −1)
2
,
Q = e−qϵB,
b(a) = 1 −2ϵ(1 + a),
β = u0 −1.
(34)
The exponential slope qϵ is the midpoint between the two contour heights u0 −1 and U −1.
The taper leaves a damping margin of order ϵ, and β = u0 −1 places the sign transition at the
target saddle. For sufficiently small ϵ, we have b > 0 on [a0, A] and B > A + 1. Define
ws(a) = −b(a)e−2a
2a2 cosh a1[a0,A](a),
wB(a) =
Q
cosh a1[B,B+1](a),
w = ws + wB.
(35)
The signed density w determines the even entire Mellin perturbation
hϵ(ζ) =
Z ∞
0
w(a)
cos(aζ) −1
 da.
(36)
For η > 0, the positive density describing the unperturbed gamma damping is
µλ,η(a) =
e−ηa
a(1 −e−2a/λ)
(a > 0).
(37)
12
Lemma 4.2. There are absolute constants ϵ0, c, C > 0 such that, for every 0 < ϵ < ϵ0, the
shells in (35) have the following properties. For every λ > 0, −1 < u ≤U, and a ∈[a0, A],
λ|ws(a)| cosh(ua) ≤(1 −cϵ)µλ,1+u(a).
At the target saddle,
Z A
a0
ws(a)a sinh(u0a) da = −1
2 log π
2 + O(ϵ),
0 ≤
Z B+1
B
wB(a)a sinh(u0a) da ≤Ce−c/ϵ2.
Finally, the target-saddle contribution and the remote-saddle domination ratio satisfy
BQe(u0−1)B ≤Ce−c/ϵ2,
C0Q−1e−(U−1)(B−A) ≤Ce−c/ϵ2.
The implicit constant in O(ϵ) is absolute.
Proof. Write
w∗(a) = −
e−2a
2a2 cosh a.
The negative shell agrees with w∗, up to its taper, on [a0, A]. Its omitted saddle contributions
are
Z a0
0
|w∗(a)|a sinh a da = O(a0),
Z ∞
A
|w∗(a)|a sinh a da = O
e−2A
A
!
,
because tanh a ≤min(a, 1). The taper changes the same integral by only
Z A
a0
|1 −b(a)| |w∗(a)|a sinh a da = O
ϵ
Z ∞
0
(1 + a)e−2a tanh a
a
da
!
= O(ϵ).
Moving from u = 1 to u = u0 contributes another O(ϵ): the mean-value theorem gives
Z A
a0
|ws(a)|a| sinh(u0a) −sinh a| da ≪ϵ
Z ∞
0
e−2a cosh(u0a)
cosh a
da ≪ϵ.
Since a0 = o(ϵ) and e−2A/A = o(ϵ), the identity
R ∞
0 w∗(a)a sinh a da = −1
2 log(π/2) now gives
the asserted negative-shell saddle displacement.
To compare the negative shell with gamma damping, divide its density by µλ,1+u:
λ|ws(a)| cosh(ua)
µλ,1+u(a)
= b(a)Θλ(a)e(u−1)a cosh(ua)
cosh a ,
Θλ(a) = 1 −e−2a/λ
2a/λ
.
The finite-dimensional correction satisfies 0 < Θλ(a) ≤1, by 1 −e−x ≤x. When −1 < u ≤1,
both e(u−1)a and cosh(ua)/ cosh a are at most 1, so the ratio is at most b(a). When 1 ≤u ≤U,
the elementary inequality cosh(ua) ≤e(u−1)a cosh a bounds it by b(a)e2(u−1)a ≤b(a)eϵa. Since
ϵA = o(1), uniformly on the negative shell,
b(a)eϵa = (1 −2ϵ(1 + a))(1 + ϵa + O(ϵ2a2)) ≤1 −cϵ.
Thus the taper retains a damping margin of order ϵ on every contour −1 < u ≤U.
It remains to check that the positive shell has opposite effects at the two saddle locations. At
u0, sinh(u0a)/ cosh a ≤e(u0−1)a, and hence
0 ≤
Z B+1
B
wB(a)a sinh(u0a) da ≤(B + 1)Qe(u0−1)(B+1).
The amplitude in (34) has exponential slope strictly between u0 −1 and U −1. Hence
Qe(u0−1)B = e−cϵB,
Qe(U−1)B = ecϵB
for an absolute c > 0. Since ϵB = ϵ−2, while B and C0 grow only polynomially in 1/ϵ, both
separation quantities are O(e−c′/ϵ2) for another absolute c′ > 0. The same estimate controls the
positive-shell saddle displacement.
□
13
The shells now determine the common envelope; it remains to impose the Fourier symmetries
and select the signs. The first gamma pole occurs at t = −iλ, or ζ = −i. Thus P0 should vanish
at −i, while P± should take the same positive value there. Reflection ζ 7→−ζ should also
exchange P−with P+. These requirements lead to the following envelope and polynomials:
Eλ(t) = πit/2Γ
λ −it
2

eλhϵ(t/λ),
P±(ζ) = 1 + ζ2 + β ± iζ(1 + ζ2),
P0(ζ) = −(1 + ζ2),
Xfj(t) = Eλ(t)Pj(t/λ),
fj(r) = r−λ
2π
Z
R
Xfj(t)rit dt
(j ∈\{−, 0, +\}).
(38)
Indeed, P−(−ζ) = P+(ζ), P0 is even, and
P±(−i) = β > 0,
P0(−i) = 0.
On the imaginary saddle branch,
P+(iu) = β + (1 −u)2(1 + u),
P−(iu) = β + (1 −u)(1 + u)2,
P0(iu) = u2 −1.
(39)
Consequently P+(iu) > 0 for every u > −1, whereas β = u0 −1 gives
P−(iu0) = ϵ
4
1 −

2 + ϵ
4
2!
< 0,
P0(iu0) =

1 + ϵ
4
2
−1 > 0.
The saddle calculation below will transfer precisely these polynomial signs to the radial func-
tions.
Lemma 4.3. For every sufficiently small ϵ > 0 and every integer d ≥1, with λ = d/2, the
inverse Mellin integrals in (38), initially defined for r > 0, extend to fj ∈Srad(Rd; R) for
j ∈\{−, 0, +\}. These extensions satisfy bf−= f+, bf0 = f0, f−(0) = f+(0) > 0, and f0(0) = 0.
Proof. Fix d and ϵ. Compact support of w makes hϵ entire, and on each horizontal line t = s+iτ
it gives
|hϵ((s + iτ)/λ)| ≤2
Z ∞
0
|w(a)| cosh(aτ/λ) da.
Thus the perturbation is bounded on every fixed horizontal strip. Uniformly for τ in compact
pole-free intervals, the standard gamma asymptotics [DLMF, § 5.11] give
|Xfj(s + iτ)| ≤Cd,ϵ,τ(1 + |s|)(λ+τ−1)/2+3e−π|s|/4.
In particular, the vertical sides of rectangular contour shifts tend to zero. Shifting the Mellin
contour upward arbitrarily far proves rapid decay as r →∞, also after differentiation. Shifting
downward past the gamma poles
t = −i(λ + 2n),
n = 0, 1, 2, . . . ,
gives an expansion in even powers r2n, with a remainder of arbitrarily high order. Thus each
fj extends to a smooth radial Schwartz function on Rd. Conjugate symmetry of its real-line
Mellin data makes this extension real.
Because hϵ is even, the envelope obeys mλ(t)Eλ(−t) = Eλ(t). The polynomial reflection
identities and the multiplier (10) therefore give
bf−= f+,
bf0 = f0.
(40)
Finally, the value at the origin is determined by the first pole of the downward-shifted contour.
More generally, Resz=−n Γ(z) = (−1)n/n! shows that the pole t = −i(λ + 2n) contributes
2πλ/2+nr2n (−1)n
n!
eλhϵ(−i(1+2n/λ))Pj(−i(1 + 2n/λ)).
(41)
14
At n = 0, evenness gives hϵ(−i) = hϵ(i), and P±(−i) = β, whereas P0(−i) = 0. Consequently,
f+(0) = f−(0) = 2πλ/2eλhϵ(i)β > 0,
f0(0) = 0.
□
(42)
4.3. Saddle geometry. Recall u0 and U from (34), and set
u∗= −1 + log λ
4λ .
(43)
To determine the signs of the Mellin integrals (38), we associate each contour t = λ(T + iu)
with the radius r = ev(u) at which Eλ(t)rit is stationary in T.
The decrease Du(T) of its
logarithmic modulus must remain positive away from T = 0. We establish this first using the
gamma contribution when u∗≤u ≤U, and then using the positive shell wB when u ≥U. The
resulting saddle estimates will identify the sign of each integral with the sign of Pj(iu).
Put
η = 1 + u > 0,
m = λη
2 ,
ψ = (log Γ)′,
ψ(1) = ψ′.
Take the branch of log Γ on Re z > 0 which is real on the positive axis.
On the contour
t = λ(T + iu), the logarithm of Eλ(t)rit is
iλ(T + iu)
2
log π + log Γ

m −iλT
2

+ λhϵ(T + iu) + iλ(T + iu) log r.
Its derivative with respect to T at T = 0 equals
iλ
1
2 log π −1
2ψ(m) −
Z ∞
0
w(a)a sinh(ua) da + log r

.
Thus T = 0 is stationary precisely when r = ev(u), where
v(u) = −1
2 log π + 1
2ψ(m) +
Z ∞
0
w(a)a sinh(ua) da,
V (u) = v′(u) = λ
4 ψ(1)(m) +
Z ∞
0
w(a)a2 cosh(ua) da.
(44)
Differentiating the saddle log-radius gives V (u) = v′(u). Thus V (u) > 0 makes ev(u) increase
with u; the quadratic decrease of the logarithmic modulus at T = 0 is λV (u)T 2/2.
To separate the gamma and shell contributions to this decrease, recall the positive density
µλ,η from (37). The standard log-gamma integral [DLMF, § 5.9], after substituting a = λs/2,
gives
Gλ,η(T) := log Γ(m −iλT/2) −log Γ(m) + iλT
2 ψ(m)
=
Z ∞
0
(eiaT −1 −iaT)µλ,η(a) da,
(45)
Dγ(T) := −Re Gλ,η(T) =
Z ∞
0
(1 −cos(aT))µλ,η(a) da.
In particular, the gamma function in the Mellin envelope always damps the integrand away
from T = 0.
Normalize Eλ(λ(T + iu))riλT by Eλ(iλu) and set r = ev(u). The linear contributions cancel
by the saddle equation, and (36) gives
Lu(T) := log Eλ(λ(T + iu))
Eλ(iλu)
+ iλTv(u)
= Gλ,η(T) + λ
Z ∞
0
w(a) cosh(ua)(cos(aT) −1) da
+ iλ
Z ∞
0
w(a) sinh(ua)(aT −sin(aT)) da,
(46)
Du(T) := −Re Lu(T)
= Dγ(T) + λ
Z ∞
0
w(a) cosh(ua)(1 −cos(aT)) da.
(47)
15
Equation (47) isolates the main difficulty: ws reduces Du(T), while wB increases it. We must
prove Du(T) > 0 for every T ̸= 0 and V (u) > 0 for every u ≥u∗. The gamma density will
dominate ws on [u∗, U], whereas wB will dominate ws on [U, ∞).
The quadratic and cubic sizes of the phase are measured by
Vγ = 1
λ
Z ∞
0
a2µλ,η(a) da = λ
4 ψ(1)(m),
M3 = 1
λ
Z ∞
0
a3µλ,η(a) da +
Z ∞
0
|ws(a)| + wB(a)
a3 cosh(ua) da.
(48)
Here Vγ is the gamma contribution to V (u), and M3 bounds the cubic error after the signed
shell contributions have been replaced by their absolute values.
For later comparisons, write
Ds(T) = λ
Z A
a0
|ws(a)| cosh(ua)(1 −cos(aT)) da,
DB(T) = λ
Z B+1
B
wB(a) cosh(ua)(1 −cos(aT)) da,
Vs =
Z A
a0
|ws(a)|a2 cosh(ua) da,
VB =
Z B+1
B
wB(a)a2 cosh(ua) da.
(49)
In particular, Du = Dγ −Ds + DB and V (u) = Vγ −Vs + VB.
Lemma 4.4. There are absolute constants c, C > 0 such that, for every λ > 0, u > −1
satisfying λ(1 + u) ≥1, and T ∈R, the phase and quantities defined in (44)–(48) satisfy
Lu(T) + λV (u)
2
T 2
 ≤CλM3|T|3.
(50)
Moreover, writing η = 1 + u,
1
2η ≤Vγ ≤C
η ,
(51)
1
λ
Z ∞
0
a3µλ,η(a) da ≤C
η2 ,
(52)
Dγ(T) ≥cλ min
T 2
η , |T|
!
(T ∈R).
(53)
Proof. The globally valid Taylor estimates
eix −1 −ix = −x2
2 + O(|x|3),
x −sin x = O(|x|3)
applied to (45) and (46), and using | sinh(ua)| ≤cosh(ua), give (50).
Since m = λη/2 ≥1/2, the variance and third-moment bounds are the standard uniform
trigamma and polygamma estimates [DLMF, §§ 5.9, 5.11, 5.15]; indeed,
Vγ = λ
4 ψ(1)(m),
1
λ
Z ∞
0
a3µλ,η(a) da = −λ2
8 ψ(2)(m).
To bound Dγ(T) at every frequency, use 1 −e−x ≤x in (37):
µλ,η(a) ≥λe−ηa
2a2 .
For T ̸= 0, take L = min(η−1, |T|−1). On 0 < a < L, both e−ηa and (1 −cos(aT))/(a2T 2) are
bounded below by absolute positive constants. Therefore Dγ(T) ≥cλT 2L, which is (53); the
case T = 0 is immediate.
□
16
The lower endpoint (43) ensures that the gamma shape parameter grows throughout the
saddle analysis:
m ≥λ(1 + u∗)
2
= 1
8 log λ.
Lemma 4.4 controls the gamma contribution and cubic remainder. We next show that, for
u∗≤u ≤U, the negative shell removes at most a (1−cϵ)-fraction of the gamma damping, while
wB contributes nonnegative damping.
Lemma 4.5. There is ϵ0 > 0 and an absolute c > 0 such that, for every 0 < ϵ < ϵ0, there are
constants Cϵ, λϵ > 0 with the following property. For every λ ≥λϵ, every u∗≤u ≤U, and
η = 1 + u,
λ|ws(a)| cosh(ua) ≤(1 −cϵ)µλ,η(a)
(a ∈[a0, A]),
(54)
Du(T) ≥cϵDγ(T)
(T ∈R),
(55)
cϵ
η ≤V (u) ≤Cϵ
η ,
(56)
M3 ≤Cϵ
η2 .
(57)
Moreover, λη ≥(log λ)/4.
Proof. Lemma 4.2 gives (54). Integrating that inequality against 1 −cos(aT) ≥0 shows that
the negative shell removes at most a (1 −cϵ)-fraction of the gamma damping. The contribution
of wB is nonnegative. Therefore
Du(T) ≥Dγ(T) −(1 −cϵ)
Z A
a0
(1 −cos(aT))µλ,η(a) da ≥cϵDγ(T),
proving (55).
For the curvature, integrating (54) against a2/λ and using (49) gives Vs ≤(1 −cϵ)Vγ, and
hence
V (u) = Vγ −Vs + VB ≥cϵVγ + VB ≥cϵ
η .
For u∗≤u ≤U, the positive-shell variance satisfies
VB ≤(B + 1)2Qe(U−1)(B+1) = Oϵ(1).
Moreover, λη ≥(log λ)/4, so (51) gives
V (u) ≤Vγ + VB = Oϵ(η−1),
because η ≤2 + ϵ/2. This proves (56).
Similarly, the negative-shell third moment is bounded by the gamma third moment:
Z A
a0
|ws(a)|a3 cosh(ua) da ≤1
λ
Z ∞
0
a3µλ,η(a) da.
The positive-shell third moment is Oϵ(1). Consequently, (52), λη ≥(log λ)/4, and η ≤2 + ϵ/2
give (57).
□
We have proved positive damping and curvature for u∗≤u ≤U. When u > U, the factor
cosh(ua) in (47) can make ws overwhelm the gamma contribution. Put δ = u −1. For T ̸= 0,
the ratios Ds(T)/(λ min(T 2, 1)) and DB(T)/(λ min(T 2, 1)) have respective sizes at most C0eδA
and at least QeδB. The separation in Lemma 4.2 therefore makes wB dominate ws throughout
u ≥U.
The interval support [B, B + 1] of wB prevents frequency resonances: a shell concentrated at
one a would have DB(T) = 0 whenever aT ∈2πZ, whereas
Z B+1
B
(1 −cos(aT)) da = 1 −sinc(T/2) cos((B + 1
2)T).
(58)
Here sinc(x) = sin(x)/x, with sinc(0) = 1. Since 1 −| sinc(T/2)| ≫min(T 2, 1), (58) is positive
for every T ̸= 0, uniformly at both small and large frequencies.
17
Define the separation error
ρϵ = (A + a−1
0 )Q−1 exp

−ϵ
2(B −A)

.
Lemma 4.2 gives ρϵ = O(e−c/ϵ2) = o(1) as ϵ ↓0. The next lemma bounds Ds(T)/DB(T) by
O(ρϵ), uniformly for u ≥U and T ̸= 0; it also gives the curvature and third-moment bounds
required for the saddle approximation.
Lemma 4.6. There are absolute constants c, C > 0 and ϵ0 > 0 such that, for every 0 < ϵ < ϵ0,
there are constants λϵ, Cϵ, cϵ > 0 with the following property. For every λ ≥λϵ, every u ≥U,
and every T ∈R, writing δ = u −1, one has
Ds(T) ≤CλC0eδA min(T 2, 1),
(59)
DB(T) ≥cλQeδB min(T 2, 1).
(60)
Consequently,
Ds(T) ≤CρϵDB(T),
(61)
Du(T) ≥Dγ(T) + cDB(T),
(62)
cVB ≤V (u) ≤CVB.
(63)
The shell variance and third moments obey
cB2QeδB ≤VB ≤(B + 1)2Qeδ(B+1),
(64)
Z A
a0
|ws(a)|a3 cosh(ua) da ≤CAρϵVB,
(65)
Z B+1
B
wB(a)a3 cosh(ua) da ≤(B + 1)VB.
In particular,
M3 ≤CϵV (u),
V (u) ≥cϵ > 0
(u ≥U).
(66)
Proof. For u = 1 + δ, the explicit shell densities satisfy
|ws(a)| cosh(ua) ≪eδa
a2 ,
wB(a) cosh(ua) ≍Qeδa.
If |T| ≤1, the bound 1 −cos(aT) = O(a2T 2) gives
Ds(T) ≪λT 2
Z A
a0
eδa da ≪λAeδAT 2.
If |T| ≥1, use 1 −cos(aT) = O(1):
Ds(T) ≪λeδA
Z A
a0
a−2 da ≪λa−1
0 eδA.
These two estimates give (59). On the positive shell, (58) yields
DB(T) ≫λQeδB
Z B+1
B
(1 −cos(aT)) da ≫λQeδB min(T 2, 1).
For T ̸= 0, division gives (61), since
Ds(T)
DB(T) ≪C0Q−1e−δ(B−A) ≤ρϵ = o(1).
At T = 0, both damping terms vanish. Comparing their quadratic coefficients in (49) gives
Vs = O(ρϵVB). Choose ϵ0 so that the implicit constant times ρϵ is less than 1/2. Then Du =
Dγ −Ds + DB proves (62), and V (u) = Vγ −Vs + VB gives
Vγ + cVB ≤V (u) ≤Vγ + VB.
Integrating wB(a) cosh(ua) ≍Qeδa against a2 gives (64). Since δ ≥ϵ/2,
VB ≫B2QeϵB/2 = B2ecϵB ≫1,
Vγ ≪η−1 ≪1.
18
Thus Vγ = O(VB), so the preceding comparison of V (u) with Vγ + VB gives (63). Since a ≤A
on the negative shell,
Z A
a0
|ws(a)|a3 cosh(ua) da ≤AVs ≪AρϵVB,
while a ≤B + 1 bounds the positive-shell third moment by (B + 1)VB. This proves (65).
Since η ≥2 + ϵ/2, (51) and (52) bound the gamma third moment by O(Vγ). The shell third
moments are Oϵ(VB), while Vγ = O(VB) and (63) gives VB ≍V (u). Hence M3 = Oϵ(V (u)).
Finally, (64) and (63) give V (u) ≫VB ≫B2QeϵB/2 =: cϵ > 0, proving (66).
□
Set T0 = (2(B + 1))−1.
Lemma 4.6 controls the total damping for u ≥U. To bound the tails of the saddle integral,
we sharpen that control on three frequency ranges: |T| ≤T0, where Du(T) is quadratic; T0 ≤
|T| ≤η, where wB supplies a uniform positive floor; and |T| ≥η, where the gamma contribution
also grows linearly.
Lemma 4.7. There is ϵ0 > 0 and an absolute c > 0 such that, for every 0 < ϵ < ϵ0, there are
constants cϵ, λϵ > 0 with the following property. For every λ ≥λϵ and u ≥U, set η = 1 + u and
δ = u −1. Then
Du(T) ≥cϵλV (u)T 2
(|T| ≤T0),
(67)
Du(T) ≥cϵλQeδB
(T0 ≤|T| ≤η),
(68)
Du(T) ≥cλ|T| + cϵλQeδB
(|T| ≥η).
(69)
Proof. First suppose that |T| ≤T0. Since |aT| ≤1/2 on [B, B + 1], 1 −cos(aT) ≫a2T 2, and
therefore
DB(T) ≫λVBT 2.
Since η ≥2 + ϵ/2 and |T| ≤η, (53) and (51) also give
Dγ(T) ≫λT 2
η
≫λVγT 2.
Combining these contributions using (62) and (63) proves (67).
Next, if T0 ≤|T| ≤η, then min(T 2, 1) ≥T 2
0 . Thus (62) and (60) give
Du(T) ≫λT 2
0 QeδB ≫ϵ λQeδB.
This proves (68).
Finally, if |T| ≥η, then min(T 2, 1) = 1 and (53) gives Dγ(T) ≫λ|T|. Adding the positive
shell through (62) and (60) yields
Du(T) ≫λ|T| + λQeδB,
which proves (69).
□
4.4. Global saddle asymptotics. The damping bounds now determine the exterior signs of
f+, f−, f0. On each contour, the centered phase is quadratic near T = 0, the factor Pj(T + iu)
is asymptotic to Pj(iu), and the remaining contour is negligible. We establish these claims
separately for the gamma-controlled range u∗≤u ≤U and the positive-shell-controlled range
u ≥U.
Lemma 4.8. Fix 0 < ϵ < ϵ0, let λ = d/2, and recall u0, U from (34) and u∗from (43). For
u > −1 and P ∈\{P+, P−, P0\}, put
Iλ,P (u) =
Z
R
eLu(T)P(T + iu) dT.
As d →∞,
Iλ,P (u) = P(iu)
s
2π
λV (u)
1 + oϵ(1)
,
(70)
19
uniformly for u ≥u∗when P = P+, and uniformly for u ≥u0 when P = P−or P = P0. More
precisely,
sup
u≥u∗
p
λV (u)Iλ,P+(u)
√
2πP+(iu)
−1
 −→0,
max
j∈\{−,0\} sup
u≥u0
p
λV (u)Iλ,Pj(u)
√
2πPj(iu)
−1
 −→0.
Proof. The poles of the integrand in (38) are t = −i(λ + 2n), n ≥0. Hence Lemma 4.3 allows
the contour to be shifted to t = λ(T + iu) whenever u > −1. At r = ev(u), the shifted Mellin
inversion formula is
fj(ev(u)) = λEλ(iλu)
2π
e−(1+u)λv(u)Iλ,Pj(u).
(71)
The prefactor of Iλ,Pj(u) is positive. By (39), P+(iu) > 0 for u > −1, while P−(iu) < 0 and
P0(iu) > 0 for u ≥u0. Their fixed degrees and uniform lower bounds on those ranges imply
|P(T + iu)|
|P(iu)|
≪ϵ 1 + |T|3,
P(T + iu)
P(iu)
= 1 + Oϵ(|T| + |T|3).
(72)
To compare Iλ,P (u) with its Gaussian approximation, choose K →∞, and set
T∗=
K
p
λV (u).
By (50), the phase on |T| ≤T∗is
Lu(T) = −λV (u)
2
T 2 + O

λM3|T|3
.
The approximation is uniform if its central interval shrinks and its cubic error tends to zero:
T∗= oϵ(1),
K3M3
√
λ V (u)3/2 = oϵ(1).
Under these conditions, (72) and the substitution x =
p
λV (u)T give
Z
|T|≤T∗
eLu(T)P(T + iu) dT = P(iu)
s
2π
λV (u)
1 + oϵ(1)
.
It remains to show that the integral over |T| > T∗is oϵ(|P(iu)|/
p
λV (u)). We verify the central
approximation and this tail bound separately on [u∗, U] and [U, ∞).
First suppose u∗≤u ≤U, and put η = 1 + u, L = λη. Equations (56) and (57) give
L ≥1
4 log λ,
V (u) ≍ϵ η−1,
M3 ≪ϵ η−2.
Choosing K = L1/12, so that K →∞and K3 = o(
√
L), gives
T∗
η ≪ϵ L−5/12,
K3M3
√
λV (u)3/2 ≪ϵ L−1/4.
The damping bounds (55) and (53) are quadratic for |T| ≤η and linear for |T| ≥η. Conse-
quently,
sup
|T|≤T∗
Lu(T) + λV (u)
2
T 2
 ≪ϵ L−1/4,
q
λV (u)
Z
T∗≤|T|≤η
(1 + |T|3)e−Du(T) dT ≪ϵ e−cϵK2,
q
λV (u)
Z
|T|≥η
(1 + |T|3)e−Du(T) dT ≪ϵ e−cϵL.
Both tail estimates tend to zero uniformly because L ≥(log λ)/4, proving (70) for u∗≤u ≤U.
20
Now suppose u ≥U, and write δ = u −1. By (63) and (66), the remote shell controls both
the curvature and the third moment:
V (u) ≍ϵ VB ≫ϵ 1,
M3 ≪ϵ V (u).
Take K = λ1/12. Then
T∗≪ϵ λ−5/12,
K3M3
√
λV (u)3/2 ≪ϵ λ−1/4.
In particular, T∗< T0 for all sufficiently large λ. The quadratic bound (67) on |T| ≤T0 yields
sup
|T|≤T∗
Lu(T) + λV (u)
2
T 2
 ≪ϵ λ−1/4,
q
λV (u)
Z
T∗≤|T|≤T0
(1 + |T|3)e−Du(T) dT ≪ϵ e−cϵK2.
(73)
For T0 ≤|T| ≤η, the variance bound (64) and the damping estimate (68) give
q
λV (u)
Z
T0≤|T|≤η
(1 + |T|3)e−Du(T) dT
≪ϵ
√
λ exp
B + 1
2
δ + 4 log(2 + δ) −cϵλQeBδ

= oϵ(1).
(74)
The exponent on the right of (74) decreases in δ ≥ϵ/2, since its derivative is (B + 1)/2 + 4/(2 +
δ) −cϵλBQeBδ < 0 for sufficiently large λ. At δ = ϵ/2, that exponent is −cϵλ + Oϵ(1). Thus
the middle-frequency contribution tends to zero uniformly even as u →∞.
Finally, for |T| ≥η, (69) supplies the positive-shell damping and a linear gamma tail, so
q
λV (u)
Z
|T|≥η
(1 + |T|3)e−Du(T) dT ≪ϵ e−cϵλ,
(75)
uniformly in δ: as in (74), the damping −cϵλQeBδ absorbs the growth of
p
V (u). Equations
(73)– (75) prove the saddle formula on every u ≥U.
□
Corollary 4.9. For every fixed 0 < ϵ < ϵ0, there is dϵ such that, for every integer d ≥dϵ,
f+(r) > 0
(r ≥ev(u∗)),
f−(r) < 0
(r ≥ev(u0)),
f0(r) > 0
(r ≥ev(u0)).
(76)
Proof. For u > −1, the prefactor of Iλ,Pj(u) in (71) is positive. Hence (70) identifies the sign
of fj(ev(u)) with that of Pj(iu). Equations (56) and (63) give v′(u) = V (u) > 0 on [u∗, ∞),
while (44) and the positive shell wB give v(u) →∞. Thus [u∗, ∞) parametrizes every radius
r ≥ev(u∗), and [u0, ∞) parametrizes every radius r ≥ev(u0). The signs in (39) now give (76).
□
4.5. Positivity and the sharp upper bound. Corollary 4.9 proves the required signs outside
the saddle radii. To finish the construction, we must also show f+(r) > 0 for 0 ≤r ≤r∗= ev(u∗).
Shifting the Mellin contour below O(log λ) gamma poles expresses f+(r)/f+(0) as a truncated
exponential series plus a uniformly negligible remainder.
Lemma 4.10. Fix 0 < ϵ < ϵ0, let λ = d/2, set r∗= ev(u∗), and write h′
1 =
R ∞
0 w(a)a sinh a da.
As d →∞,
sup
0≤r≤r∗
eπe2h′
1r2 f+(r)
f+(0) −1
 −→0.
In particular, f+(r) > 0 on [0, r∗] for all sufficiently large d.
Proof. First identify the range on which the truncated exponential series must approximate e−y.
Put
y = πe2h′
1r2,
H(u) =
Z ∞
0
w(a)a sinh(ua) da,
η∗= 1 + u∗= log λ
4λ .
21
The normalized contour height u∗tends to −1, the normalized height of the first gamma pole,
while the gamma shape parameter λ(1 + u∗)/2 = (log λ)/8 still diverges. This choice makes
(70) applicable and keeps y(r∗) logarithmic. Indeed, H is odd and, for fixed ϵ, has bounded
derivative near −1. Consequently
H(u∗) + h′
1 = H(−1 + η∗) −H(−1) = Oϵ(η∗).
By (44),
y(r∗) = exp

ψ
λη∗
2

+ 2H(u∗) + 2h′
1

.
The standard digamma asymptotic [DLMF, § 5.11] and λη∗/2 = (log λ)/8 therefore give
0 ≤y ≤y(r∗) = 1
8 log λ + Oϵ(1).
(77)
Since e−y can be as small as a negative power of λ on (77), the approximation error must be
controlled relative to e−y.
To recover enough exponential-series coefficients, set
N = ⌈log λ⌉,
p = N + 1
2,
κ = 1 + 2p
λ .
Since p = N + 1/2, the contour t = s −i(λ + 2p) lies strictly between consecutive gamma poles.
Furthermore, p ≍log λ makes the exponential-series tail negligible on (77). To justify shifting
to this contour, we verify that the multiplier eλhϵ(t/λ) preserves a uniform exponential decay
margin. Indeed, pA/λ = oϵ(1), so the negative-shell taper satisfies
b(a)e2pa/λ ≤1 −cϵ
(a0 ≤a ≤A)
for an absolute c > 0 and all sufficiently large λ. For every intermediate contour height 0 ≤q ≤κ,
the positive shell contributes nonpositively to Re hϵ(s/λ −iq) −hϵ(iq), and the negative shell
gives
λ
Re hϵ(s/λ −iq) −hϵ(iq)
 ≤(1 −cϵ)λ
2
Z ∞
0
1 −cos(as/λ)
a2
da
= (1 −cϵ)π|s|
4 .
Standard gamma asymptotics [DLMF, § 5.11] supply the complementary factor e−π|s|/4. Hence
the integrand decays like e−cϵ|s| on the vertical sides, permitting the shift in (38) to t = s −
i(λ + 2p). This shift crosses exactly the poles t = −i(λ + 2n), 0 ≤n ≤N.
Applying the residue formula (41) and the origin value (42) gives
f+(r)
f+(0) =
N
X
n=0
(−y)n
n!
Aλ,n + Rλ,p(r),
(78)
Aλ,n = eλ[hϵ(i(1+2n/λ))−hϵ(i)]−2nh′
1 P+(−i(1 + 2n/λ))
β
,
|Aλ,n −1| ≪ϵ
n(1 + n)
λ
(0 ≤n ≤N).
(79)
Indeed, Taylor expansion at u = 1 gives
λ

hϵ

i

1 + 2n
λ

−hϵ(i)

= 2nh′
1 + Oϵ
n2
λ
!
,
P+(−i(1 + 2n/λ))
β
= 1 + Oϵ
n
λ

,
uniformly for n ≤N; here N2/λ = o(1).
For r > 0, the remainder Rλ,p(r) in (78) is the integral over t = s −i(λ + 2p):
Rλ,p(r) = πλ/2+pr2p
2πf+(0)
Z
R
πis/2Γ(−p −is/2)
· eλhϵ(s/λ−iκ)P+(s/λ −iκ)ris ds.
22
The gamma reflection and product estimates [DLMF, §§ 5.5, 5.8], together with p ∈Z + 1
2, give
the standard bound
|Γ(−p −is/2)| ≪e−π|s|/4
Γ(1 + p),
λ
 Re hϵ(s/λ −iκ) −hϵ(iκ)
 ≤(1 −cϵ)π|s|
4 .
Since P+(s/λ −iκ)/β = Oϵ(1 + |s|3), integration in s now gives
|Rλ,p(r)| ≪ϵ
yp
Γ(1 + p).
(80)
Here we used
λ[hϵ(iκ) −hϵ(i)] = 2ph′
1 + Oϵ(p2/λ),
(πr2)pe2ph′
1 = yp,
and absorbed p2/λ = o(1). The estimate extends to r = 0 by continuity.
To compare (78) with e−y, we bound the coefficient errors Aλ,n −1, the contour remainder
Rλ,p(r), and the omitted terms P
n>N(−y)n/n!. Equations (79), (80), and (77) give
ey
N
X
n=0
yn
n! |Aλ,n −1| ≪ϵ
(1 + y)2e2y
λ
≪ϵ
(log λ)2
λ3/4
,
ey|Rλ,p(r)| ≪ϵ
eyyp
Γ(1 + p),
ey X
n>N
yn
n! ≪eyyN+1
(N + 1)!.
(81)
The first estimate follows from P
n≥0 n(1 + n)yn/n! = (y2 + 2y)ey.
For the two tails, put
L = log λ. Since y ≤L/8 + Oϵ(1) and p = L + O(1), Stirling’s formula [DLMF, § 5.11] gives
log
eyyp
Γ(1 + p) ≤
1
8 + 1 −log 8

L + Oϵ(log L).
The same estimate holds with p replaced by N + 1. Since 1/8 + 1 −log 8 < 0, both tails are
smaller than the coefficient error in (81). Therefore (78) yields, uniformly on 0 ≤r ≤r∗,
f+(r)
f+(0) = e−y
1 + Oϵ
(log λ)2
λ3/4
!!
> 0
(0 ≤r ≤r∗).
□
(82)
Proof of Theorem 4.1. Set Rϵ,d = ev(u0). The Fourier identities and origin values follow from
(40) and (42). Corollary 4.9 gives the required exterior signs, while (82) supplies positivity of
f+ on the remaining interval. Thus
bf−= f+ > 0,
f−(0) = f+(0) > 0,
bf0 = f0,
f0(0) = 0,
f−(r) < 0 < f0(r)
(r ≥Rϵ,d).
(83)
For fixed ϵ, the saddle equation (44) and the digamma asymptotic [DLMF, § 5.11] give
lim
d→∞
Rϵ,d
√
d
=
r1 + u0
4π
exp
Z ∞
0
w(a)a sinh(u0a) da

.
(84)
The two shell contributions in Lemma 4.2 give
Z ∞
0
w(a)a sinh(u0a) da = −1
2 log π
2 + O(ϵ).
Since u0 →1, (33) and (84) give
lim
ϵ↓0 lim
d→∞
Rϵ,d
√
d
= 1
π.
□
(85)
23
Proof of Theorem 1.1. By (83), the dilation Fϵ,d(x) = f−(Rϵ,dx) is admissible and satisfies
Fϵ,d(0)/ bFϵ,d(0) = Rd
ϵ,d by Fourier scaling. Inserting Fϵ,d into (3) gives
LPd ≤vd
2d Rd
ϵ,d.
(86)
The universal lower bound (28), the admissible upper bound (86), and (85), together with
v1/d
d
√
d →
√
2πe, imply
r e
2π ≤lim inf
d→∞LP1/d
d
≤lim sup
d→∞
LP1/d
d
≤
√
2πe
2π
=
r e
2π.
□
Proof of Theorem 1.2. The lower bound is Proposition 3.7. For the upper bound, define gϵ,d,−=
f+ −f−and gϵ,d,+ = f0. Equation (40) and Fourier inversion give
bgϵ,d,−= −gϵ,d,−,
bgϵ,d,+ = gϵ,d,+.
By (83), both gϵ,d,−and gϵ,d,+ vanish at the origin and are strictly positive outside Rϵ,d; in par-
ticular, neither is zero. The definition (6) therefore gives Aς(d) ≤Rϵ,d for both signs. Combining
this bound with (85) gives
1
π ≤lim inf
d→∞
Aς(d)
√
d
≤lim sup
d→∞
Aς(d)
√
d
≤lim
ϵ↓0 lim
d→∞
Rϵ,d
√
d
= 1
π.
□
Appendix A. Comparison of the sign-uncertainty constants
For an anti-self-Fourier radial function g, the central Mellin moment Mg(d/2) vanishes. In-
tegrating the radial tail of g therefore produces a self-Fourier function with a strictly smaller
last-sign radius.
Proposition A.1. Let d ≥1, put λ = d/2, and suppose that 0 ̸= g ∈L1
rad(Rd; R) satisfies
bg = −g and g(0) = 0. Define Tdg(0) = 0 and, for x ̸= 0, set
(Tdg)(x) = λ
2
Z ∞
1
tλ−1g(tx) dt,
(87)
where g denotes its continuous Fourier-inversion representative. Then Tdg is nonzero, continu-
ous, and integrable, with
d
Tdg = Tdg,
(Tdg)(0) = 0,
∥Tdg∥1 ≤1
2∥g∥1.
If r(g) < ∞, then r(Tdg) < r(g); if g is Schwartz, then so is Tdg. Consequently, the constants
in (6) satisfy
A+(d) < A−(d)
(d ≥1).
Proof. Since g = −bg ∈L1, Fourier inversion makes g bounded and continuous. In particular,
R
Rd |g(x)| |x|−λ dx < ∞, by splitting the integral at |x| = 1. The Mellin–Fourier identity in (9)
was established for Schwartz functions, so we first verify its central consequence directly for
the integrable eigenfunction g. Set J(t) =
R
Rd g(x)e−πt|x|2 dx. Gaussian duality and Tonelli’s
theorem give
J(t) = −t−λJ(1/t),
Z ∞
0
tλ/2−1|J(t)| dt ≤Γ(λ/2)
πλ/2
Z
Rd |g(x)| |x|−λ dx < ∞.
Hence the substitution t 7→1/t makes
R ∞
0 tλ/2−1J(t) dt equal to its negative. Gaussian integra-
tion and polar coordinates therefore give the central Mellin cancellation
Mg(λ) =
Z ∞
0
g(r)rλ−1 dr = 0.
(88)
The integral defining Tdg(x) converges absolutely for x ̸= 0, since g is integrable and sλ−1 ≪x
sd−1 for s ≥|x|. Tonelli’s theorem and d = 2λ yield ∥Tdg∥1 ≤(λ/2)∥g∥1
R ∞
1 t−λ−1 dt = ∥g∥1/2,
24
also justifying Fourier transformation under the integral. Fourier scaling and (88) give
d
Tdg(ξ) = λ
2
Z ∞
1
tλ−d−1bg(ξ/t) dt
= −λ
2
Z 1
0
sλ−1g(sξ) ds
= λ
2
Z ∞
1
sλ−1g(sξ) ds = Tdg(ξ).
(89)
The small-scale representation in (89) extends continuously to x = 0, with value −g(0)/2 = 0.
Differentiating the large-scale representation gives (x · ∇+ λ)Tdg = −λg/2, so Tdg ̸= 0. If g
is Schwartz, differentiating the small-scale representation gives smoothness at the origin, and
differentiating the large-scale representation gives rapid decay at infinity; thus Tdg is Schwartz.
If R = r(g) < ∞, then R > 0, since otherwise g ≥0 and
R g = bg(0) = −g(0) = 0 would force
g = 0. Hence g(s) ≥0 for s ≥R, and (5) and (87) give
(Tdg)(r) = λ
2 r−λ
Z ∞
r
sλ−1g(s) ds > 0
(r ≥R).
Indeed, equality at any r ≥R would make both g and bg = −g compactly supported, con-
tradicting Fourier analyticity.
In particular, Tdg(R) > 0, so continuity gives r(Tdg) < R.
Applying this strict decrease to an attained radial negative-sign extremizer [CG19, Thm. 1.4]
gives A+(d) < A−(d).
□
References
[AJCHLT20] N. Afkhami-Jeddi, H. Cohn, T. Hartman, D. de Laat, and A. Tajdini, High-dimensional sphere
packing and the modular bootstrap, J. High Energy Phys. 12 (2020), article 066, doi:10.1007/
JHEP12(2020)066.
[Ahl79]
L. V. Ahlfors, Complex Analysis, 3rd ed., McGraw–Hill, New York, 1979.
[BCK10]
J. Bourgain, L. Clozel, and J.-P. Kahane, Principe d’Heisenberg et fonctions positives, Ann. Inst.
Fourier (Grenoble) 60 (2010), 1215–1232, doi:10.5802/aif.2552.
[Coh02]
H. Cohn, New upper bounds on sphere packings. II, Geom. Topol. 6 (2002), 329–353, doi:10.2140/
gt.2002.6.329.
[CDG24]
H. Cohn, D. Dong, and F. Gonçalves, Sign uncertainty principles and low-degree polynomials, Proc.
Amer. Math. Soc. Ser. B 11 (2024), 224–228, doi:10.1090/bproc/219.
[CE03]
H. Cohn and N. Elkies, New upper bounds on sphere packings. I, Ann. of Math. (2) 157 (2003),
689–714, doi:10.4007/annals.2003.157.689.
[CG19]
H. Cohn and F. Gonçalves, An optimal uncertainty principle in twelve dimensions via modular
forms, Invent. Math. 217 (2019), 799–831, doi:10.1007/s00222-019-00875-4.
[CKMRV17] H. C
\end{document}
